\subsection{Round seven: research log}
\label{sec:r7-appendix}
\paragraph{2026-10-26, round seven.} Peter (2026-10-26), verbatim: ``In the Norns questions, I really think we should have a bunch of VTR placements that generate data that we then analyze with the qeustion what might be better clustering and intitialization - every idea we've had is in the other direction - hypothesis, test, not really better - it's like my data mining Journal paper with Gharibian, Shannon and Wilton''. I take the paper he means to be Gharibian, Shannon, Jamieson and Chung, TRETS 2013 \cite{Gharibian2013tret}, with the FPL 2014 paper by Gharibian, Shannon and Jamieson \cite{gharibian2014identifying}; this identification is an assumption of mine (his list names Wilton, and the Crossref entries of both papers, checked on 2026-10-26, do not list an author of that name). Question N Q1.7 (area, speed, wirelength): how far is the VPR packer from a best-known clustering made with hindsight from good placements, and how much would LUT duplication (after Schabas and Brown \cite{schabas2003duplication}) add? Question N Q2.4: which HDL relations predict which atoms sit together in good placements, and where does stock packing disagree? The answers are Tables~\ref{tab:r7gap} to~\ref{tab:r7coef}. The protocol text was committed before the first run (\code{docs/plan/round7.md}, commit \code{93c983a}); the run code was not (it was committed as \code{94827f5} after the first run), so the commit froze the text only. The remote records the creation of the branch with that commit before the first run finished (GitHub event log); the protocol's timestamp entries give the details.

\begin{table}[t]
\caption{Research questions added in round seven.}
\label{tab:r7rq}
\centering\scriptsize
\begin{tabular}{@{}lp{6.4cm}@{}}
\toprule
RQ & question (measures) \\
\midrule
N Q1.7 & On real designs, how far is the VPR packer from a best-known clustering made with hindsight from good placements, and how much would LUT duplication add? (area, speed, wirelength) \\
N Q2.4 & Mined from many placements linked to the HDL: which HDL relations (word, bit index, statement, instance) predict which atoms sit together in good placements, and where does stock packing disagree? (wirelength, speed) \\
\bottomrule
\end{tabular}
\end{table}

\paragraph{Agent defaults.} The hindsight clustering is written as a \code{.net} (logic elements as units, exact CLB legality of the architecture) and run with packing off, from the best run of the previous round, arm chosen and evaluated on seeds 1 to 5 (in sample); the cost weight $1+4\,\mathrm{pcrit}$, the move radius of two CLBs are fixed in the protocol; carry-chain CLBs stay as they are; the affinity-bias arm of the first batch is kept and renamed. The best affinity-bias run is used for the disagreement table and the figures; the \code{.net} clustering is used as a mining target. The raw VPR directories are kept in scratch directories (15.7~GB for the first batch) and \code{runs/r7} holds 0.8~GB; nothing was deleted.

Regeneration: \code{python scripts/write\_round7.py} from \code{data/r7/*.json}; the steps that need \code{runs/} are \code{r7.py hind}, \code{r7net.py run}, \code{r7corpus.py}, \code{r7stats.py}, \code{r7fstats.py} and \code{r7mine.py}.
